\phantomsection
\chapter*{译者序}
\markboth{译者序}{译者序}
\addcontentsline{toc}{chapter}{译者序}
本书是对布朗大学Sergei Treil教授的著作《Linear Algebra Done Wrong》（更新于2024年10月1日）的翻译。
全书共分为9章，具体内容见后文“前言”一节，建议读者仔细阅读之，调整好阅读本书的心态和目标。

{\it 待完成：本书内容概要}

% 本书第一章引入一系列基本概念，首先是向量空间和线性组合，然后先引入基，再分别从存在性和唯一性引出张成和线性无关的概念。接下来介绍线性变换，借助矩阵的语言表述有关$\mathbb{F}^n$的情况，注重阐述了矩阵乘法的动机。然后介绍了可逆和同构。最后还给出了一小节应用实例。第二章。
% 第三章从低维几何角度（体积）引入行列式，同时也给出了基于置换的正式定义
% 本书对预备知识的要求低，语言通俗易懂，论述详细，
% 因此可作为非数学类学生的进阶自学用书，教师也可借此教材加开第二学期线性代数进阶课程。

% 特别感谢原书作者Sergei Treil教授及其同事的大力支持与热情鼓励。
本翻译使用的模板是在 \href{https://github.com/andy123t}{andy123t} 编写的 \href{https://github.com/andy123t/zhbook}{zhbook} 基础上修改而成的，
感谢 \href{https://github.com/andy123t}{andy123t} 精心打造这一美观的模板。

英文的表达习惯与中文有诸多不同，且原书语言生动，若直译难免生硬别扭、失其神韵。因此，
译者在充分尊重原文句意的前提下，对其中部分字眼不求字字对应，而根据自己的理解作了补充，必要时在脚注中作出解释。
翻译本质上是一种再创作，无法代替原著。希望读者不仅阅读翻译，更要能独立阅读原著，或者至少对照阅读一遍，
想必会更有收获。

由于译者只能利用业余时间进行翻译且水平有限，本翻译一定有不妥之处甚至错误，欢迎读者们多提意见和建议。
% 请将意见和建议发送至 \mailto{ladr2024@163.com}。

原书基于CC-BY-NC-ND 3.0许可协议发布，按照规定是不允许再创作的，此翻译版只是译者怀着帮助读者跨越“语言关”、更深入地掌握书中内容的朴素想法所作，译者不希望因此招致麻烦。所以，希望各位仅在课程群内交流，不要外传。谢谢大家。

{\hfill 译者\hspace{2em}}

{\hfill 2026年3月}